\uuid{j88h}
\titre{Maximum de variables exponentielles et loi de Gumbel}

\niveau{L3}
\module{Modélisation de l'aléatoire}
\chapitre{Convergence en loi et valeurs extrêmes}
\sousChapitre{Maximum de variables exponentielles}
\theme{Loi de Gumbel et modélisation des crues}
\auteur{Erwan Hillion}
\datecreate{2026-10-01}
\organisation{ESM3/ISD}
\difficulte{4}

\contenu{

\texte{
Soit $(X_i)_{i\geq 1}$ une famille de variables aléatoires mutuellement indépendantes, suivant chacune la loi exponentielle $\mathcal E(\lambda)$, avec $\lambda>0$. On pose, pour tout $n\geq 1$,
$$
M_n=\max(X_1,\ldots,X_n).
$$
On définit également
$$
Y_n=\lambda M_n-\ln n.
$$
Enfin, on considère une variable aléatoire $Z$ de densité
$$
g(x)=e^{-x}e^{-e^{-x}},\qquad x\in\mathbb R.
$$
}

\begin{enumerate}

\item \question{Donner la fonction de répartition $F$ de $X_1$.}
\reponse{
Comme $X_1\sim\mathcal E(\lambda)$, sa densité est
$$
f_{X_1}(x)=\lambda e^{-\lambda x}\mathbf 1_{[0,+\infty[}(x).
$$
Ainsi,
$$
F(x)=\mathbb P(X_1\leq x)
=
\begin{cases}
0,&x<0,\\[1mm]
1-e^{-\lambda x},&x\geq 0.
\end{cases}
$$
}

\item \question{Montrer que la fonction de répartition de $M_n$ est donnée par
$$
F_{M_n}(x)=\mathbb P(M_n\leq x)=F(x)^n.
$$}
\reponse{
On a
$$
\{M_n\leq x\}
=
\{X_1\leq x\}\cap\cdots\cap\{X_n\leq x\}.
$$
Les variables $X_1,\ldots,X_n$ étant indépendantes,
$$
\mathbb P(M_n\leq x)
=
\prod_{k=1}^n\mathbb P(X_k\leq x)
=F(x)^n.
$$
Par conséquent,
$$
F_{M_n}(x)
=
\begin{cases}
0,&x<0,\\[1mm]
\left(1-e^{-\lambda x}\right)^n,&x\geq 0.
\end{cases}
$$
}

\item \question{Montrer que les fonctions de répartition de $Y_n$ et $M_n$ sont liées par
$$
F_{Y_n}(x)=F_{M_n}\left(\frac{x+\ln n}{\lambda}\right).
$$}
\reponse{
Pour tout $x\in\mathbb R$,
\begin{align*}
F_{Y_n}(x)
&=\mathbb P(Y_n\leq x)\\
&=\mathbb P(\lambda M_n-\ln n\leq x)\\
&=\mathbb P\left(M_n\leq\frac{x+\ln n}{\lambda}\right).
\end{align*}
Donc
$$
\boxed{F_{Y_n}(x)=F_{M_n}\left(\frac{x+\ln n}{\lambda}\right).}
$$
}

\item \question{Montrer que, pour tout $x>0$,
$$
F_{Y_n}(x)=\left(1-\frac{e^{-x}}{n}\right)^n.
$$}
\reponse{
Pour $x>0$, on a $x+\ln n>0$. En utilisant la question précédente,
\begin{align*}
F_{Y_n}(x)
&=\left(1-\exp\left[-\lambda\frac{x+\ln n}{\lambda}\right]\right)^n\\
&=\left(1-e^{-x-\ln n}\right)^n\\
&=\left(1-\frac{e^{-x}}{n}\right)^n.
\end{align*}
En réalité, cette formule est valable dès que $x\geq-\ln n$. Si $x<-\ln n$, alors $F_{Y_n}(x)=0$.
}

\item \question{Montrer que la fonction de répartition de $Z$ est
$$
F_Z(x)=e^{-e^{-x}},\qquad x\in\mathbb R.
$$}
\reponse{
La fonction
$$
G(x)=e^{-e^{-x}}
$$
vérifie
$$
G'(x)=e^{-x}e^{-e^{-x}}=g(x).
$$
De plus,
$$
\lim_{x\to-\infty}G(x)=0
\qquad\text{et}\qquad
\lim_{x\to+\infty}G(x)=1.
$$
Ainsi $G$ est bien la fonction de répartition associée à la densité $g$, et
$$
\boxed{F_Z(x)=e^{-e^{-x}}.}
$$
}

\item \question{Montrer que la suite $(Y_n)_{n\geq1}$ converge en loi vers $Z$.}
\reponse{
Fixons $x\in\mathbb R$. Pour $n$ suffisamment grand, on a $x+\ln n>0$, donc
$$
F_{Y_n}(x)=\left(1-\frac{e^{-x}}{n}\right)^n.
$$
Or, pour tout $a\geq0$,
$$
\left(1-\frac{a}{n}\right)^n\longrightarrow e^{-a}.
$$
En prenant $a=e^{-x}$, on obtient
$$
F_{Y_n}(x)\longrightarrow e^{-e^{-x}}=F_Z(x).
$$
La fonction $F_Z$ étant continue sur $\mathbb R$, le théorème de caractérisation de la convergence en loi donne
$$
\boxed{Y_n\xrightarrow{\mathcal L}Z.}
$$
}

\item \question{On modélise la hauteur maximale d'un fleuve pendant l'année $i$ par une variable $X_i\sim\mathcal E(0{,}5)$, les années étant indépendantes. Interpréter $M_{100}$.}
\reponse{
La variable
$$
M_{100}=\max(X_1,\ldots,X_{100})
$$
représente la plus grande hauteur maximale annuelle observée au cours d'une période de cent ans.
}

\item \question{On approxime la loi de $Y_{100}$ par celle de $Z$. Calculer $\mathbb P(Z>a)$ avec
$$
a=-\ln\bigl(-\ln(0{,}99)\bigr).
$$}
\reponse{
Par définition de $a$,
$$
e^{-a}=-\ln(0{,}99).
$$
Ainsi,
$$
F_Z(a)=e^{-e^{-a}}
=e^{\ln(0{,}99)}
=0{,}99.
$$
Donc
$$
\boxed{\mathbb P(Z>a)=1-F_Z(a)=0{,}01.}
$$
}

\item \question{On donne $a\simeq4{,}6$ et $\ln(100)\simeq4{,}6$. Déterminer $b>0$ tel que
$$
\mathbb P(M_{100}>b)\simeq0{,}01.
$$}
\reponse{
Comme
$$
Y_{100}=\lambda M_{100}-\ln(100),
$$
on choisit $b$ de sorte que
$$
\lambda b-\ln(100)=a.
$$
Alors
$$
\mathbb P(M_{100}>b)
=\mathbb P(Y_{100}>a)
\simeq\mathbb P(Z>a)
=0{,}01.
$$
Avec $\lambda=0{,}5\ \mathrm m^{-1}$,
$$
b=\frac{a+\ln(100)}{\lambda}
\simeq\frac{4{,}6+4{,}6}{0{,}5}
=18{,}4.
$$
Ainsi,
$$
\boxed{b\simeq18{,}4\ \mathrm m.}
$$
}

\end{enumerate}

}
